\section{Why filter after chaining but before alignment?}
\label{sec:prefilter-role}

Seed matches and chains identify sequence intervals that may align well. They
do not establish that the complete intervals satisfy an edit-distance bound.
Repeated sequence, short seeds, hash collisions, and sequencing errors can all
produce plausible chains whose intervening gaps are expensive to align and
ultimately fail the requested bound. Applying affine-gap dynamic programming
to every candidate is correct, but it spends time on candidates that a cheaper
procedure may be able to disprove.

A pre-alignment filter addresses this case. It runs after chaining has defined
candidate regions and before the base aligner processes chain-defined fills.
Its result is deliberately limited. The filter may show that a candidate
cannot meet the requested threshold, or it may leave the question unresolved.
It does not produce an alignment, a CIGAR string, or an edit script.

This distinction is required by the SQL interface. Consider a predicate of the
form

\begin{verbatim}
WHERE mmsearch(sequence, query)
  AND bioseq.edit(sequence) < 10
\end{verbatim}

The strict bound \(B=10\) means that the unit-cost edit count reconstructed
from every returned alignment traceback must be below ten. The prefilter
receives the equivalent inclusive threshold

\[
    E = B - 1 = 9.
\]

For mandatory fill \(i\), let \(D_i\) be the minimum global Levenshtein
distance between its query and target strings. Let \(C_i\) be the edit count
on the traceback selected by the base aligner. A safe filter computes \(L_i\)
such that

\[
    0 \leq L_i \leq D_i \leq C_i.
\]

The second inequality holds because no particular edit script can use fewer
edits than the minimum over all scripts. For disjoint mandatory fills,
\[
    \sum_i L_i \leq \sum_i C_i \leq C_{\mathrm{region}},
\]
where \(C_{\mathrm{region}}\) is the edit count exposed by
\texttt{bioseq.edit}. A region may therefore be rejected when
\(\sum_iL_i>E\). A retained candidate still proceeds to the selected aligner,
which constructs a traceback before the original strict predicate is applied.
The approximation can remove work, but it cannot replace exact evaluation.

The terms false positive and false negative are easy to misread in this
setting. Here, the positive filter outcome means that a candidate is retained
for alignment. The two errors have different effects.

\begin{table}[htbp]
\centering
\small
\begin{tabular}{
    >{\raggedright\arraybackslash}p{0.22\textwidth}
    >{\raggedright\arraybackslash}p{0.31\textwidth}
    >{\raggedright\arraybackslash}p{0.36\textwidth}}
\textbf{Outcome} & \textbf{Meaning} & \textbf{Effect}\\
\hline
False positive & A pair above the threshold is retained & Additional alignment
work, with no change to the query result\\
False negative & A pair within the threshold is rejected & A valid result is
lost, which violates the SQL predicate\\
\end{tabular}
\caption{Consequences of prefilter errors. Safe filtering permits false
positives but not false negatives.}
\label{tab:prefilter-errors}
\end{table}

BioSeqDB applies this contract to the geometry created by chaining. For a
multi-seed chain, mandatory fill windows extend from the first anchor start to
successive anchor ends. These windows include the ranks covered by anchors.
An exact or hash-supported seed is not assigned zero edit cost merely because
it formed an anchor. Its symbols are compared within the mandatory fill, so a
hash collision can loosen candidate generation but cannot remove edits from
the filter bound. Single-seed spans are handled separately by exact edit
accounting.

The terminal windows have different semantics. Local extension and X-drop may
omit an expensive prefix or suffix. A large global distance in such a window
does not prove that the permitted local alignment fails. The prefilter
therefore examines mandatory fill windows and leaves optional end extensions
to the aligner. Reverse-strand fills are first expressed in the orientation in
which they will be aligned.

Unequal lengths also require a restricted rule. Every transformation from
\(Q\) to \(T\) needs at least

\[
    \Delta_{\mathrm{len}} = \bigl|\lvert Q\rvert-\lvert T\rvert\bigr|
\]

insertions or deletions. BioSeqDB uses this exact lower bound for unequal
lengths rather than extending an equal-length approximation beyond its stated
domain. Unsupported sizes, thresholds, and execution geometries are retained
conservatively. Resource or device failures are ordinary query errors, since
silently changing the selected execution route would make performance and
diagnostic behavior difficult to interpret.

\section{SneakySnake chip-maze filtering}
\label{sec:prefilter-sneakysnake}

SneakySnake treats approximate matching as routing through a chip maze
\cite{sneakysnake_paper}. For equal-length strings of length \(m\) and an
inclusive threshold \(E\), the maze has \(2E+1\) rows. Row \(d\), where
\(d\in[-E,E]\), represents one displacement between the query and target. At
target position \(j\), its cell is

\[
 Z_d(j)=
 \begin{cases}
 0 & 0\leq j+d<m \text{ and } Q[j+d]=T[j]\\
 1 & \text{otherwise.}
 \end{cases}
\]

A zero is a match on that diagonal and a one is an obstacle. Starting at a
checkpoint, the algorithm measures the matching prefix on every available
diagonal. It follows the longest run, crosses the obstacle at its end when the
right boundary has not yet been reached, and starts another search from the
new checkpoint. The number of crossed obstacles is the filter value. If it
exceeds \(E\), the pair cannot satisfy the threshold in the safe equal-length
domain.

Figure~\ref{fig:prefilter-chip-maze} shows three escape segments. The
highlighted route stays on the unshifted row. It follows exact runs at
positions 0--2, 4--6, and 8--10, crossing obstacles at positions 3, 7, and 11.
The third crossing raises the count above \(E=2\), so this example is rejected.
The maze abstracts away the edit script. In particular, it records neither the
operation responsible for a change of row nor the vertical distance between
selected segments.

\begin{figure}[H]
\centering
\begingroup
\small
\setlength{\tabcolsep}{4pt}
\renewcommand{\arraystretch}{1.2}
\newcommand{\mazeMatch}{\cellcolor{white}\(\cdot\)}
\newcommand{\mazeObstacle}{\cellcolor{black!70}\textcolor{white}{\(\times\)}}
\newcommand{\mazeRoute}{\cellcolor{orange!35}\textcolor{orange!80!black}{\(\rightarrow\)}}
\newcommand{\mazeCross}{\cellcolor{orange!70}\(\times\)}
\begin{tabular}{c|*{12}{c}}
 & 0 & 1 & 2 & 3 & 4 & 5 & 6 & 7 & 8 & 9 & 10 & 11\\
\hline
\(d=+2\) & \mazeMatch & \mazeMatch & \mazeMatch & \mazeObstacle &
\mazeMatch & \mazeMatch & \mazeMatch & \mazeObstacle & \mazeMatch &
\mazeMatch & \mazeObstacle & \mazeObstacle\\
\(d=+1\) & \mazeMatch & \mazeMatch & \mazeMatch & \mazeObstacle &
\mazeMatch & \mazeMatch & \mazeMatch & \mazeObstacle & \mazeMatch &
\mazeMatch & \mazeMatch & \mazeObstacle\\
\(d=0\) & \mazeRoute & \mazeRoute & \mazeRoute & \mazeCross &
\mazeRoute & \mazeRoute & \mazeRoute & \mazeCross & \mazeRoute &
\mazeRoute & \mazeRoute & \mazeCross\\
\(d=-1\) & \mazeObstacle & \mazeMatch & \mazeMatch & \mazeObstacle &
\mazeMatch & \mazeMatch & \mazeMatch & \mazeObstacle & \mazeMatch &
\mazeMatch & \mazeMatch & \mazeObstacle\\
\(d=-2\) & \mazeObstacle & \mazeObstacle & \mazeMatch & \mazeObstacle &
\mazeMatch & \mazeMatch & \mazeMatch & \mazeObstacle & \mazeMatch &
\mazeMatch & \mazeMatch & \mazeObstacle\\
\end{tabular}

\vspace{0.4em}
\begin{tabular}{clclcl}
\cellcolor{white}\(\cdot\) & match &
\cellcolor{black!70}\textcolor{white}{\(\times\)} & obstacle &
\cellcolor{orange!35}\(\rightarrow\) & selected escape segment
\end{tabular}
\endgroup
\caption{A rejecting SneakySnake chip-maze walk for
\(Q=\mathtt{AAAAAAAAAAAA}\), \(T=\mathtt{AAACAAACAAAC}\), and \(E=2\).
Each row represents a query-to-target displacement. Amber cells show the
successive escape segments and three crossed obstacles on the unshifted row.}
\label{fig:prefilter-chip-maze}
\end{figure}

The corresponding procedure is compact.

\begin{algorithm}[H]
\caption{Unpartitioned SneakySnake core for one equal-length submaze}
\label{alg:prefilter-sneakysnake}
\small
\begin{algorithmic}[1]
\State \(c\gets0\)
\State \(r\gets0\)
\State \textbf{while} \(c<m\) \textbf{do}
    \State \(b\gets\max_{d\in[-E,E]}
        \operatorname{matchingPrefix}(Q,T,c,d)\)
    \State \textbf{if} \(c+b\geq m\) \textbf{then return} retain
    \State \(r\gets r+1\)
    \State \textbf{if} \(r>E\) \textbf{then return} reject
    \State \(c\gets c+b+1\)
\State \textbf{end while}
\State \textbf{return} retain
\end{algorithmic}
\end{algorithm}

The lower-bound direction follows from the relation between a maze route and
an edit script. Between two edits, an alignment follows one diagonal through a
run of equal symbols. A substitution ends a run on the same diagonal. An
insertion or deletion moves the next run to another diagonal. An edit script
of cost \(k\) therefore induces matching segments separated by at most \(k\)
obstacles. The chip maze also permits transitions whose cost an edit script
would represent more precisely, so it is a relaxation of edit distance.

SneakySnake chooses the furthest reachable checkpoint under this relaxed
model. The greedy escape rule minimizes the number of obstacles in the routing
model~\cite{sneakysnake_paper}. If \(D(Q,T)\leq E\), the edit script induces a
maze route with at most \(E\) obstacles, and the greedy route cannot use more.
Consequently, a greedy count above \(E\) proves that the global distance also
exceeds \(E\). The converse does not hold because the relaxation can omit
costs that an edit script must pay.

The maze is less expressive than edit-distance dynamic programming. This loss
of information is intentional. The obstacle count can be smaller than the
true distance, which retains some pairs that later fail alignment. It must not
be larger in the domain used for rejection. Partitioning a long gap into
bounded submazes may weaken the result further, since a new submaze forgets
some preceding routing state. That change can add false positives, but it does
not create false negatives. Bounds from non-overlapping mandatory pieces can
still be accumulated because those pieces require separate edits.

BioSeqDB compares normalized alphabet ranks rather than DNA characters. The
same equality operation can therefore handle nucleotide and amino-acid
alphabets once both operands have a common representation. A scalar CPU path
scans matching runs directly. A vector CPU path compares contiguous groups of
ranks and locates the first mismatch. Both execute the same checkpoint rule.
On a GPU, separate candidates can be assigned to independent invocations at
small thresholds. Larger thresholds expose more diagonal work, so a small
cooperating group can evaluate shifts for one candidate. Fixed submaze sizes
bound memory use and keep dispatch geometry predictable.

Long gaps are divided into consecutive bounded submazes, and
Algorithm~\ref{alg:prefilter-sneakysnake} is applied to each one. The
lower-bound counts are added. Restarting at a submaze boundary discards routing
state and can reduce the resulting count. It cannot increase the bound above
the value for the complete maze, so partitioning weakens filtering rather than
changing accepted query results.

At each checkpoint, an exact pruning rule can reduce prefix scans. Suppose the
best run found so far has length \(b\). Another diagonal can improve the
maximum only if its first \(b\) positions match, including the position at
offset \(b\). Testing that position first discards only a diagonal that cannot
win. Every possible improvement still receives a complete matching-prefix
comparison.

These mappings change how comparisons are scheduled, not the rejection
criterion. If a vector lane, work group, or bounded representation cannot
express a job, the job is retained for exact alignment. This rule keeps
backend selection separate from query correctness.

\section{GateKeeper shifted-mask filtering}
\label{sec:prefilter-gatekeeper}

SneakySnake must repeatedly search diagonals at each checkpoint. Very short
gaps with small thresholds admit a more regular bit-vector computation.
GateKeeper derives such a filter from Shifted Hamming Distance
methods~\cite{gatekeeper_paper,gatekeeper_gpu_paper}. For equal-length strings,
the unshifted mismatch mask is

\[
    H_0[j]=[Q[j]\neq T[j]].
\]

For each displacement \(d=1,\ldots,E\), two further masks compare the target
with shifted query positions.

\[
\begin{aligned}
H_{-d}[j] &= [j\geq d \land Q[j-d]\neq T[j]]\\
H_{+d}[j] &= [j+d<m \land Q[j+d]\neq T[j]].
\end{aligned}
\]

The complete representation contains \(2E+1\) masks. A zero denotes a possible
match. A run of true matches displaced from the main diagonal can therefore
appear in one of the shifted masks even when \(H_0\) contains mismatches.

Isolated zeros are often accidental matches rather than evidence of a useful
run. GateKeeper amends each mask before the masks are combined. BioSeqDB fills
an internal zero run of length one or two only when ones bound it on both
sides. Longer runs and runs touching an edge remain unchanged. For an
unbounded bit vector \(M\), with positions outside the valid interval treated
as zero, this operation can be written as

\[
\begin{aligned}
S &= \neg M \land (M\ll1) \land (M\gg1)\\
P &= \neg M \land (\neg M\gg1) \land (M\ll1) \land (M\gg2)\\
\operatorname{amend}(M) &= M\lor S\lor P\lor(P\ll1).
\end{aligned}
\]

The amended result is restricted to the valid sequence positions. When a mask
spans several machine words, neighboring bits must be carried across word
boundaries. Otherwise an identical two-zero run would receive different
treatment depending on its position in memory.

A shift also creates a vacant position at one end. If this position remained
zero, the later intersection could interpret it as a match and hide an end
error. After amendment, BioSeqDB sets the boundary-adjacent bits
\[
    \widehat H_{-d}[d-1]\gets1,
    \qquad
    \widehat H_{+d}[m-d]\gets1.
\]
The index identifies the position next to the vacant prefix or suffix for
displacement \(d\). The ordering is significant because compensation before
amendment could turn an edge run into a bounded run and fill it.

All amended and compensated masks are intersected. The result is divided into
four-bit groups. A zero group contributes no cost. The patterns

\[
0101,\quad 0110,\quad 1001,\quad 1010,\quad 1011,\quad 1101
\]

contribute two, while every other nonzero group contributes one. Accumulation
can stop at \(E+1\), since larger values have the same threshold outcome. This
grouped estimate is not a popcount and it is not an exact edit distance.

\begin{figure}[H]
\centering
\small
\renewcommand{\arraystretch}{1.15}
\begin{tabular}{c}
\fbox{\parbox{0.74\textwidth}{\centering
Construct the unshifted mismatch mask \(H_0\)}}\\[0.35em]
\(\operatorname{popcount}(H_0)\leq E\): retain \qquad
\(E=0\): reject by exact Hamming comparison\\[0.35em]
\(\Downarrow\) otherwise\\[0.35em]
\fbox{\parbox{0.74\textwidth}{\centering
Construct the \(2E\) shifted masks, amend each mask independently,\\
then set the boundary-adjacent compensation bits}}\\[0.35em]
\(\Downarrow\)\\[0.35em]
\fbox{\parbox{0.74\textwidth}{\centering
Intersect the masks and compute the capped four-bit group estimate}}\\[0.35em]
\(\Downarrow\) retain when the estimate is at most \(E\)\\[0.35em]
\fbox{\parbox{0.74\textwidth}{\centering
For a proposed rejection, evaluate exact thresholded edit distance}}\\[0.35em]
\(D(Q,T)\leq E\): retain \qquad \(D(Q,T)>E\): reject
\end{tabular}
\caption{GateKeeper decision flow in BioSeqDB. The unshifted mask provides an
exact rejection certificate when \(E=0\). For positive thresholds, shifted
masks can retain a candidate or propose rejection, and the exact threshold
guard confirms the proposal.}
\label{fig:prefilter-gatekeeper}
\end{figure}

\subsection{The exact rejection guard}

Mask amendment changes zeros to ones, and the grouped estimate is an
approximation. An observed absence of false negatives on a data set would not
prove safety for every sequence. BioSeqDB therefore treats an estimate above
\(E\) as a rejection proposal. A thresholded exact Levenshtein calculation
confirms whether \(D(Q,T)>E\). If it does not, the candidate is retained.

The exact guard changes the role of the published filter. GateKeeper was
designed to predict whether a pair is likely to satisfy an edit threshold, and
its filtering accuracy is evaluated empirically
~\cite{gatekeeper_paper,gatekeeper_gpu_paper}. An empirical false-negative
rate is not a sufficient contract for a database predicate, where one rejected
valid pair changes the result set. BioSeqDB therefore uses shifted masks to
decide whether exact threshold evaluation is needed. The masks can retain a
candidate immediately, but for \(E>0\) they cannot reject one without the
guard. When \(E=0\), a nonzero unshifted mismatch mask is itself an exact
certificate because no insertion or deletion can produce a zero-edit
alignment.

\begin{algorithm}[H]
\caption{Guarded GateKeeper decision for one equal-length gap}
\label{alg:prefilter-gatekeeper}
\small
\begin{algorithmic}[1]
\State \(H\gets\operatorname{unshiftedMismatchMask}(Q,T)\)
\State \textbf{if} \(\operatorname{popcount}(H)\leq E\)
    \textbf{then return} retain
\State \textbf{if} \(E=0\) \textbf{then return} reject
\State \(A\gets\operatorname{amend}(H)\)
\State \textbf{for} \(d\gets1\) to \(E\) \textbf{do}
    \State \(H_{-d}\gets\operatorname{shiftedMismatchMask}(Q,T,-d)\)
    \State \(H_{+d}\gets\operatorname{shiftedMismatchMask}(Q,T,+d)\)
    \State \(M_-\gets\operatorname{amend}(H_{-d})\)
    \State \(M_-[d-1]\gets1\)
    \State \(M_+\gets\operatorname{amend}(H_{+d})\)
    \State \(M_+[m-d]\gets1\)
    \State \(A\gets A\land M_-\land M_+\)
\State \textbf{end for}
\State \textbf{if} \(\operatorname{groupedCost}(A)\leq E\)
    \textbf{then return} retain
\State \textbf{if} \(\operatorname{exactDistanceExceeds}(Q,T,E)\)
    \textbf{then return} reject
\State \textbf{return} retain
\end{algorithmic}
\end{algorithm}

The exact guard evaluates only the threshold question. The scalar and generic
GPU routes use banded dynamic programming. A path that leaves the band already
requires more than \(E\) insertions or deletions. Values are capped at \(E+1\),
and only the preceding and current rows are needed. Specialized GPU variants
instead use an exact bounded wavefront with longest-common-prefix extension.
Both forms answer whether \(D(Q,T)>E\), use bounded memory, and omit traceback.
They run only when the mask calculation proposes rejection.

This design changes how GateKeeper results may be combined. A retained mask
estimate is not a proved lower bound, so it must not be added across fills. A
confirmed rejection is sufficient on its own because one mandatory fill already
exceeds the region threshold. The mask remains useful as a cheap screening
step that limits calls to the exact guard.

The rank representation again removes assumptions about nucleotide encoding.
On CPUs, masks fit into one or a small number of machine words for the intended
short-gap domain. Boolean operations then evaluate many sequence positions at
once. On GPUs, each invocation handles a candidate with regular register
operations. Specialized kernels are useful for common short lengths and low
thresholds, while a generic kernel covers other supported mask geometries.
Longer gaps or larger thresholds are routed to SneakySnake or retained for the
base aligner.

\section{Dispatch and query-pipeline integration}
\label{sec:prefilter-dispatch}

Algorithm selection depends on both candidate geometry and execution backend.
The scalar and vector CPU routes use SneakySnake across their supported range.
A GPU batch uses GateKeeper when all jobs are short, equal in length, and have
small thresholds. Other supported GPU batches use SneakySnake. Resolving the
algorithm once for a batch avoids separate transfers and result merging. A
mixed batch that does not satisfy GateKeeper's domain follows the broader
SneakySnake route.

The mapping is summarized in Table~\ref{tab:prefilter-routing}. The table
describes policy rather than query semantics. Every route must satisfy the
same rejection contract.

\begin{table}[htbp]
\centering
\small
\begin{tabular}{
    >{\raggedright\arraybackslash}p{0.22\textwidth}
    >{\raggedright\arraybackslash}p{0.28\textwidth}
    >{\raggedright\arraybackslash}p{0.38\textwidth}}
\textbf{Execution route} & \textbf{Preferred filter} & \textbf{Reason}\\
\hline
Scalar CPU & SneakySnake & Direct prefix scans have low setup cost and cover
general supported gaps\\
Vector CPU & SneakySnake & Contiguous rank comparisons accelerate matching
prefix detection\\
GPU, short regular gaps & Guarded GateKeeper & Bit masks map to regular
per-candidate operations\\
GPU, broader supported gaps & SneakySnake & Parallel diagonal evaluation
covers larger thresholds without large mask state\\
Unsupported geometry & No rejection & Exact alignment preserves query
semantics\\
\end{tabular}
\caption{High-level mapping of prefilter algorithms to execution routes.}
\label{tab:prefilter-routing}
\end{table}

The relevant stage order is

\[
\begin{gathered}
\text{seeding}
\rightarrow \text{chaining}
\rightarrow \text{preliminary classification}\\
\rightarrow \text{mandatory-fill prefiltering}\\
\rightarrow \text{base alignment}
\rightarrow \text{exact edit predicate}
\rightarrow \text{final result classification}
\end{gathered}
\]

This order has two important consequences. First, rejected regions do not
consume base-alignment work. Retained regions are unchanged, including scoring
parameters, gap costs, local-extension rules, and traceback. Second, the exact
predicate runs before final primary and secondary classification. A region that
violates the SQL bound cannot alter the final classification or mapping quality
of a region that satisfies it.

Each fill result is associated with its owning region. SneakySnake lower bounds
for disjoint mandatory fills are added with saturation at \(E+1\). A region is
removed when that total exceeds \(E\). A GateKeeper mask estimate is not
composable and contributes no numeric lower bound. GateKeeper rejects a region
only through the exact \(E=0\) Hamming certificate or an exact-guard
confirmation. Survivor compaction retains region order so downstream
classification sees the same ordering it would have seen without
prefiltering.

The prefilter is activated only when the query asks for exact edit accounting
and supplies a usable positive strict bound. Bounds that cannot be represented
are not converted by unsigned wraparound. The exact predicate still handles
zero, negative, and otherwise unsupported cases according to SQL semantics.
Backend choice follows the query's execution plan. Merely having a GPU
available does not move an otherwise CPU query to a device.

\section{Limitations}
\label{sec:prefilter-limitations}

\paragraph{Retention is not evidence of similarity.}
SneakySnake may undercount edits because it omits vertical movement between
escape segments and resets state at bounded submazes. GateKeeper may combine
zeros from different shifted masks. Both effects can retain a pair whose true
distance exceeds \(E\). Exact alignment and the final predicate remain
necessary.

\paragraph{Only mandatory fills are filtered.}
The method does not reject from a costly local prefix or suffix because the
aligner may legally omit it. This restriction reduces the amount of work that
can be filtered, but it preserves local-extension semantics.

\paragraph{Unequal lengths receive a weaker bound.}
The length difference can reject an impossible threshold, but it says little
about substitutions within unequal strings. Such candidates often proceed to
alignment even when their true distance is large.

\paragraph{Algorithm domains are finite.}
Bit masks, bounded submazes, and GPU work groups impose length and threshold
limits. Inputs outside a proved and representable domain are retained. A new
kernel or encoding cannot be used for rejection until it satisfies the same
lower-bound or exact-guard contract.

\paragraph{Filtering has a break-even point.}
A filter adds rank comparisons, batching, and dispatch overhead. If nearly all
candidates survive, this cost may exceed the saved alignment work. GPU launch
latency is particularly relevant for small batches. Dispatch policy therefore
uses candidate geometry, but performance selection never weakens the exact
post-alignment check.

\paragraph{The edit model is specific.}
The filter bounds minimum unit-cost Levenshtein distance, while the SQL
predicate counts unit-cost edits on the selected traceback. The latter count
can be larger, which is the safe direction established by
\(D_i\leq C_i\). The aligner may select that traceback using affine gap scores
or other scoring choices. A bound for another distance model must not be
treated as a proof for this predicate.

\section{Conclusion}

BioSeqDB places prefiltering between chaining and base alignment, where it can
remove regions whose mandatory fills cannot satisfy an exact SQL edit bound.
SneakySnake
provides a general lower-bound route through shifted matching diagonals.
GateKeeper provides a compact shifted-mask calculation for short regular gaps,
with an exact threshold guard before rejection. CPU and GPU implementations
schedule these calculations differently while retaining the same contract.
Unsupported inputs and inconclusive estimates proceed to alignment, and every
survivor is checked by the exact predicate. This asymmetry permits extra work
from false positives while preventing result loss from false negatives.
